\section{Harness Engineering：工具、接口与控制流共设计}

\subsection{为什么 harness 是一等公民}
Sutton 在 00:26 明确说出一句本讲高频引用的话：\textit{``The design of the agent harness and the tools, you have to co-design it with whatever the model can do.''}  
这句话意味着：我们不该把模型当黑盒再外挂工具，而应把工具接口、提示格式、控制流结构当成与模型同级的设计对象。

\begin{importantbox}{Harness Engineering 的三条工程原则}
\begin{enumerate}
    \item 动作语义要贴近模型熟悉分布，降低调用歧义；
    \item 工具粒度要足够紧凑，减少 ``为做一件事调用五个工具''；
    \item 输出格式要可压缩且可追踪，避免上下文被噪声淹没。
\end{enumerate}
\end{importantbox}

\subsection{Agent Computer Interface（ACI）视角}
课堂中引用了 ``ACI''（Agent Computer Interface）概念，意在类比 HCI：人机界面设计有范式，智能体与执行环境的接口同样需要系统方法。比如 ``跳转到定义''、``按符号检索引用'' 这类动作，往往比通用 grep 更符合模型思维路径，也更接近 IDE 工作流。

\begin{knowledgebox}{工具设计的“熟悉路径”原则}
主讲人借用了类似 ``Little Red Riding Hood principle'' 的说法：尽量把模型引导到后训练中见过的大量交互范式。工具名、参数、返回文本越贴近常见开发语境，越容易稳定调用成功。
\end{knowledgebox}

\begin{figure}[H]
\centering
\includegraphics[width=0.88\textwidth]{frame-03-harness.jpg}
\caption{课程中强调 ``harness 与工具必须联合设计'' 的关键画面 \protect\footnotemark}
\end{figure}
\footnotetext{视频画面时间区间：00:26:24--00:26:55。}

\subsection{控制流设计谱系：从自由到约束}
讲座把 coding agent 控制流放在一条连续谱上：一端是 ``给 shell + IDE，几乎全动态''；中间是 SWE-agent 一类 ReAct 循环；再往右是分阶段拼接（先信息检索、后补丁生成）；最右是显式状态机（state machine）限制每一步可执行动作。

\begin{table}[H]
\centering
\begin{tabular}{p{3.2cm}p{4.4cm}p{3.8cm}p{2.6cm}}
\toprule
\textbf{控制流范式} & \textbf{优点} & \textbf{风险} & \textbf{适配任务} \\
\midrule
全动态 ReAct & 灵活探索、适应未知情况 & 轨迹漂移、成本不稳定 & 开放式定位与调试 \\
两阶段拼接循环 & 信息收集与补丁生成解耦 & 阶段边界错了会级联失败 & 中等复杂修复任务 \\
状态机约束循环 & 可控性强、便于审计与复现 & 搜索空间受限、可能错过解 & 高风险场景、合规流程 \\
\bottomrule
\end{tabular}
\caption{控制流设计是任务风险与探索能力之间的权衡}
\end{table}

\begin{warningbox}{“控制越强越好”是误解}
过强约束会让 Agent 失去在异常路径中的恢复能力；过弱约束又会产生长轨迹幻觉。正确做法是按任务复杂度与风险等级分层配置控制流，而不是追求单一范式。
\end{warningbox}

\subsection{轨迹分析：必须看失败样本}
在安全 agent 讨论前，Sutton 重复了一个机器学习老原则：\textit{``Always look at the data.''} 这里的 ``data'' 不是训练集，而是 Agent 执行轨迹。只有逐条查看失败路径，才能知道是检索错位、工具误用、还是验证器误导。

\begin{importantbox}{实操建议：轨迹驱动的 harness 迭代}
\begin{itemize}
    \item 每次版本迭代固定抽样失败轨迹做人工审阅；
    \item 将高频失败模式映射到新工具或新约束；
    \item 把 ``修完一次'' 的经验沉淀到 system instruction 与策略模板中。
\end{itemize}
\end{importantbox}

\begin{figure}[H]
\centering
\includegraphics[width=0.88\textwidth]{frame-04-control-loop.jpg}
\caption{围绕控制流与策略循环的讲解，强调动态循环和程序化约束并非对立 \protect\footnotemark}
\end{figure}
\footnotetext{视频画面时间区间：00:38:10--00:38:55。}

\subsection{本章小结}
Harness Engineering 的核心是 ``系统协同'' 而非 ``Prompt 微调''。工具语义、反馈结构、控制流约束和模型能力必须联合优化，才可能在复杂任务中稳定提升。

